\begin{abstract}
K-Waay原协议已要求同次\batchreceive 输入对应不同参与方。本文构造共享生命周期的两槽Tamarin模型，受控比较放宽、消息级与参与方级配置。结果表明：放宽配置可重复接受同一精确发送来源；消息级配置虽满足消息区分和作用域内发生注入性，同一参与方见证仍可达；参与方级正控制保持参与方区分，且有效批次可达。结论仅适用于两槽批处理接纳抽象。

\noindent\textbf{关键词：}K-Waay；批处理接纳；不同参与方条件；Tamarin；形式化验证
\end{abstract}

\begin{center}
  \bfseries Abstract
\end{center}
\noindent
The complete K-Waay specification already requires the inputs of one
\batchreceive invocation to correspond to different parties.  To clarify the
same-batch relation maintained by this condition and to test whether controls
over message or sender-occurrence coordinates can replace it, we construct
three fixed-two-slot Tamarin models with a common sender-origin and sequential
processing lifecycle.  The controlled comparison covers relaxed admission, a
message-level restriction, and a party-level positive control.  The relaxed
configuration admits a trace in which one exact sender origin supports two
receiver events in the same batch context.  The message-level configuration
satisfies message distinction and scoped occurrence injectivity, while two
different messages from one party can still be accepted together.  The
party-level configuration preserves same-batch party distinction and retains
a reachable valid distinct-party batch.  Hence, message distinction and
exact-origin occurrence injectivity do not substitute for the required party
relation.  The conclusion is limited to the fixed-two-slot admission
abstraction and does not constitute an attack on complete K-Waay executions
or a theorem for arbitrary batch sizes.

\noindent\textbf{Key words:} K-Waay; batch admission; distinct-party
condition; Tamarin; formal verification

